% !TeX root = ../navier-stokes-manuscript.tex
% ============================================================
% 2.1 Vortex Stretching
% ============================================================

Axial stretching lengthens one narrow material vortex-tube segment; because incompressibility contracts its cross-section, rotating fluid moves inward and the interior spins faster. This is the positive stretching contribution. Whether the rotation strengthens overall still depends on viscosity and the rest of the signed evolution. The velocity field carrying this segment is governed by the momentum equation:
\[
  \underbrace{\partial_t u}_{\text{local acceleration}}
  +
  \underbrace{(u\cdot\nabla)u}_{\text{nonlinear transport}}
  =
  \underbrace{-\nabla p}_{\text{pressure}}
  +
  \underbrace{\nu\Delta u}_{\text{viscous diffusion}}.
\]
Taking its curl turns nonlinear transport into vorticity advection and stretching, removes the pressure gradient, and gives the material rotational balance used below.

\subsection{Vortex Stretching}\label{sub:ThePerfectlyStretchingVortex}

Keep the same segment in a locally axisymmetric region, with its vorticity aligned to the stretching axis and the strain effectively uniform across its narrow width. If \(e\) denotes that axis and \(L(t)\) the segment's length, then
\[
  S:=\frac12\bigl(\nabla u+(\nabla u)^{\mathsf T}\bigr),
  \qquad
  Se=\sigma(t)e,
  \qquad
  \sigma(t)>0,
  \qquad
  \frac{\dot L(t)}{L(t)}
  =
  e\cdot Se
  =
  \sigma(t).
\]
Its length grows at the positive fractional rate \(\sigma(t)\). Its material volume \(\mathcal V\) cannot change, so its effective transverse area is \(A(t):=\mathcal V/L(t)\), and its effective radius is \(r_{\mathrm{eff}}(t):=\sqrt{A(t)/\pi}\). Incompressibility then gives
\[
  \nabla\cdot u=0,
  \qquad
  \frac{d}{dt}(A L)=0,
  \qquad
  \frac{\dot A}{A}=-\sigma(t),
  \qquad
  \frac{\dot r_{\mathrm{eff}}}{r_{\mathrm{eff}}}
  =
  -\frac{\sigma(t)}{2}.
\]
As the segment lengthens, its area contracts at rate \(\sigma(t)\) and its effective radius contracts at half that rate. For \(\omega:=\nabla\times u\), the curl of the momentum equation gives the vorticity carried by this same segment:
\[
  \frac{D\omega}{Dt}
  =
  S\omega+\nu\Delta\omega,
  \qquad
  \frac{D}{Dt}
  :=
  \partial_t+u\cdot\nabla,
\]
and the local alignment \(\omega=|\omega|e\) converts its stretching part into
\[
  S\omega
  =
  \sigma(t)\omega,
  \qquad
  \frac12\frac{D|\omega|^2}{Dt}
  =
  \sigma(t)|\omega|^2
  +
  \nu\,\omega\cdot\Delta\omega.
\]
The first term is positive and captures the spin-up caused by contraction toward the axis. Net rotational growth occurs only where that positive contribution exceeds the signed viscous contribution. The energy identity proved in \Cref{prop:ClassicalKineticEnergyIdentity}, together with the antisymmetric part of \(\nabla u\), gives
\[
  \norm{u(t)}_2
  \leq
  \norm{u_0}_2
  <\infty,
  \qquad
  \left|\frac12\bigl(\nabla u-(\nabla u)^{\mathsf T}\bigr)\right|_{\mathrm F}
  =
  \frac{|\omega|}{\sqrt2}
  \leq
  |\nabla u|_{\mathrm F}.
\]
Finite kinetic energy can thus coexist with concentration of the rotational part of the velocity gradient inside the narrowing segment.

\divider{\NavierStokes{} Scaling}

That coexistence calls for a full-field comparison at a finer width. At a time \(t_0<T_*\), set \(\ell:=r_{\mathrm{eff}}(t_0)\). For \(\lambda>1\), define
\[
  u_\lambda(x,t)
  :=
  \lambda u(\lambda x,\lambda^2t),
  \qquad
  p_\lambda(x,t)
  :=
  \lambda^2p(\lambda x,\lambda^2t).
\]
The field \(u_\lambda\) is another solution, and at time \(\lambda^{-2}t_0\) its corresponding segment has effective radius \(\lambda^{-1}\ell\). It is a comparison at another scale, not a later material state of the opening segment. Each part of its momentum evolution changes by the same power:
\[
\begin{aligned}
  \partial_tu_\lambda(x,t)
  &=
  \lambda^3(\partial_tu)(\lambda x,\lambda^2t),
  \\
  (u_\lambda\cdot\nabla)u_\lambda(x,t)
  &=
  \lambda^3\bigl((u\cdot\nabla)u\bigr)(\lambda x,\lambda^2t),
  \\
  \nabla p_\lambda(x,t)
  &=
  \lambda^3(\nabla p)(\lambda x,\lambda^2t),
  \\
  \nu\Delta u_\lambda(x,t)
  &=
  \lambda^3\nu(\Delta u)(\lambda x,\lambda^2t).
\end{aligned}
\]
At the finer width, nonlinear transport, viscosity, pressure, and local acceleration all acquire the factor \(\lambda^3\). Viscosity gains no relative advantage from the change of scale.

\divider{Critical Height}

Although the momentum balance remains matched, homogeneous Sobolev norms change according to their order. Let \(\Lambda:=(-\Delta)^{1/2}\). The Fourier transform of the rescaled field is
\[
  \widehat{u_\lambda}(\xi,t)
  =
  \lambda^{-2}\widehat u(\lambda^{-1}\xi,\lambda^2t),
\]
and changing variables in the homogeneous norm gives
\[
  \norm{u_\lambda(t)}_{\dot H^s}^2
  =
  \lambda^{2s-1}
  \norm{u(\lambda^2t)}_{\dot H^s}^2.
\]
Norms below order \(\tfrac12\) lose powers of \(\lambda\), norms above it gain powers, and the exponent vanishes at \(s=\tfrac12\). Define the critical height by
\[
  H(t)
  :=
  \frac12\norm{u(t)}_{\dot H^{1/2}}^2
  =
  \frac12\norm{\Lambda^{1/2}u(t)}_2^2.
\]
For the rescaled solution, write \(H_\lambda\) for the same half-derivative quantity. The kinetic-energy norm, critical height, and first-derivative norm then satisfy
\[
  \norm{u_\lambda(t)}_2^2
  =
  \lambda^{-1}\norm{u(\lambda^2t)}_2^2,
  \qquad
  H_\lambda(t)=H(\lambda^2t),
  \qquad
  \norm{\nabla u_\lambda(t)}_2^2
  =
  \lambda\norm{\nabla u(\lambda^2t)}_2^2.
\]
The rescaled field has less kinetic energy, the same critical height, and a larger first-derivative norm at corresponding times. Scale neutrality at the middle order fixes no sign for \(H'(t)\).

\divider{Nonlinear Production versus Viscous Dissipation}

The sign comes from nonlinear production against viscous dissipation. At critical height, define the signed production by
\[
  \mathcal P_{1/2}[u](t)
  :=
  -\operatorname{Re}
  \pair
    {\Lambda^{1/2}u(t)}
    {\Lambda^{1/2}\bigl((u(t)\cdot\nabla)u(t)\bigr)}_{L^2}.
\]
Because \(\nabla\cdot u=0\), pressure contributes no pairing, while viscosity contributes \(-\nu\norm{\Lambda^{3/2}u}_2^2\). Consequently,
\[
  H'(t)
  +
  \nu\norm{\Lambda^{3/2}u(t)}_2^2
  =
  \mathcal P_{1/2}[u](t).
\]
The critical height rises exactly when signed nonlinear production exceeds positive viscous dissipation. Under rescaling, those two quantities change together:
\[
  \mathcal P_{1/2}[u_\lambda](t)
  =
  \lambda^2\mathcal P_{1/2}[u](\lambda^2t),
  \qquad
  \nu\norm{\Lambda^{3/2}u_\lambda(t)}_2^2
  =
  \lambda^2\nu\norm{\Lambda^{3/2}u(\lambda^2t)}_2^2.
\]
Both acquire the same factor \(\lambda^2\), so changing scale still gives neither one relative ground. The corresponding physical interval between nonlinear scales has not been determined.

\divider{The Heat Clock}

A linear viscous comparison supplies a time associated with width, without asserting the duration of a nonlinear transition. For the heat equation \(v_t=\nu\Delta v\), Fourier modes evolve by
\[
  \widehat v(\xi,t+\delta t)
  =
  e^{-\nu|\xi|^2\delta t}\widehat v(\xi,t).
\]
A scale-localized velocity structure of width \(r\) is concentrated near frequencies \(|\xi|\simeq r^{-1}\). Its heat evolution changes by order one once
\[
  \nu|\xi|^2\delta t\simeq1,
\]
so its linear viscous comparison time is
\[
  \tau_\nu(r)
  :=
  \frac{r^2}{\nu}.
\]
Halving the width quarters this comparison time. At a general finer width,
\[
  \tau_\nu(\lambda^{-1}r)
  =
  \lambda^{-2}\tau_\nu(r).
\]
The linear heat comparison contracts time by the same scale factor as the full \NavierStokes{} rescaling, but it does not set the actual lifetime of a nonlinear structure.

\divider{The Zeno Cascade}

Now sample dyadic widths. For
\[
  r_n:=2^{-n}r_0,
  \qquad
  \tau_n:=\frac{r_n^2}{\nu},
\]
the associated heat times obey
\[
  \tau_n
  =
  \frac{r_0^2}{\nu}4^{-n},
  \qquad
  \sum_{n=0}^{\infty}\tau_n
  =
  \frac{4r_0^2}{3\nu}
  <
  \infty.
\]
Their finite sum is only an available schedule. To realize it, one \NavierStokes{} velocity field would have to exhibit scale-localized structures at widths
\[
  r_0>r_1>r_2>\cdots\longrightarrow0
\]
and at sampled times
\[
  t_0<t_1<t_2<\cdots\uparrow T_*<\infty,
  \qquad
  t_{n+1}-t_n\asymp\tau_n,
\]
with comparison constants independent of \(n\). Under that condition, the time remaining after the \(n\)-th sample satisfies
\[
  T_*-t_n
  \asymp
  \sum_{k=n}^{\infty}\tau_k
  =
  \frac{4r_n^2}{3\nu}.
\]
The next sampled gap and the total time left are both comparable to \(r_n^2/\nu\), and both collapse as the sampled widths approach zero. Even this conditional one-field history does not control the first or successive Eulerian time derivatives of that velocity field at one fixed spatial point.